\scanpage{29}
\mspage{27}
% archive sheet number: 28
% hand: Grothendieck
% sheet_kind: cancelled_draft
% element: text | cancelled top, lines 1–3 | 4.10 On peut aussi regarder
\deleted{\unsure{4.10} On peut aussi regarder, \deleted{\unsure{pour $P$}} pour une représentation \deleted{\unsure{$E$}} $\ell$-adique donnée de $I$, le sous-module}
% element: display | cancelled top, line 4 | E^P
\[
\deleted{E^P}
\tag*{(4.10.4)}
\]
% element: text | cancelled upper, lines 5–7 | des invariants sous P
\deleted{des invariants sous $P$. Il est stable sous $I$, et par passage au quotient $I/P=I_t\simeq\prod_{\ell\ne p}\mathbb Z_\ell(1)$ opère sur $E^P$.}

% element: text | upper-middle, lines 8–10 | En tous cas
En tous cas, une représentation $\ell$-adique quasi-unipotente de $\pi$ sera déterminée par la donnée de $U$, et d'une représentation \uline{algébrique} de $\Gamma_{\mathbb Q_\ell}\times\unsure{W|\mathbb W}(\mathbb Q_\ell(1))$.
% element: margin note | right boxed margin, beside section 6 | On trouve le système projectif
\margin{On trouve \unsure{le système|les systèmes} projectif de \unsure{gpes|petits} Alg sur $\mathbb Q_\ell$ $\Gamma_i\times W(\mathbb Q_\ell(1))$ ($\Gamma_i$ quotients $I/V_i$, $V_i$ ss-grps ouverts invariants) dont les ``repr. alg'' donnent une \uline{$\otimes$-catégorie équivalente à celle des repr. $\ell$-adiques de $I$ qui sont quasi-unipotentes}.}

% element: text | middle, section title | 6. Cas d'un corps résiduel fini
6. \uline{Cas \deleted{\unsure{de}} d'un corps résiduel \inserted{$k$} fini}.
% element: margin note | left margin, slanted, beside section 6 | généraliser à un corps résiduel
% slanted
\margin{\deleted{\unsure{généraliser} à un corps résiduel de \unsure{type} fini, plus gén. tel que l'ext $k'\subset\overline k$ de $k$ obtenue en adjoignant les racines $\ell^\nu$-ièmes de $1$ (tous $\nu$) soit infinie.}}
% element: text | cancelled middle, line 11 | Si
\deleted{Si \deleted{\unsure{$k$}}}
% element: display | cancelled middle, line 12 | Card k=q=p^f
\[
\deleted{\operatorname{card}k=q=p^f}
\tag*{(6.1)}
\]
% element: text | cancelled middle, line 13 | on sait que
\deleted{on sait que}
% element: display | cancelled middle, line 14 | pi_0 isom Zhat
\[
\deleted{\pi_0\simeq\widehat{\mathbb Z},}
\tag*{(6.2)}
\]
% element: text | cancelled middle, lines 15–16 | l'isomorphisme étant défini
\deleted{\deleted{\unsure{Alors}} l'isomorphisme étant défini par l'élément $\operatorname{frob}_q$ de $\pi_0$, défini par}
% element: display | cancelled middle-lower, line 17 | frob_q(lambda)=lambda^q
\[
\deleted{\operatorname{frob}_q(\lambda)=\lambda^q\qquad(\lambda\in\overline k).}
\tag*{(6.3)}
\]
% element: text | cancelled lower, lines 18–20 | Cet automorphisme
\deleted{Cet automorphisme opère sur $\overline k^*$ par élévation à la puissance $q$, donc il opère de même sur les $\mu_{\unsure{n|\nu}}(\overline k^*)$, donc sur $I_t=\prod_{\ell\ne p}\mathbb Z_\ell(1)$.}
% element: text | cancelled lower, lines 21–22 | Par suite
\deleted{Par suite, \unsure{une|si} \unsure{ou une|repr.} représentation \deleted{$\ell$-adique}\inserted{continue} donnée (sur un corps \deleted{topologique} quelconque $R$)}
% element: display | cancelled lower, line 23 | u:pi to Aut_R(E)
\[
\deleted{u:\pi\unsure{-|\longrightarrow}\operatorname{Aut}_R(E)}
\]
% element: text | cancelled bottom, lines 24–26 | est triviale
\deleted{est \deleted{nulle sur} triviale sur le sous-groupe $\unsure{I^0|I'|I^p}$ de $I$ \deleted{\unsure{engendré}}, alors, considérant $I_t$ comme isomorphe à un quotient de}
% element: display | cancelled bottom, line 27 | Zhat isom product Z_l
\[
\deleted{\widehat{\mathbb Z}\simeq\prod_\ell\mathbb Z_\ell}
\]
